\section{Impacto socioeconómico y distribución de la riqueza}

\subsection{Un choque más severo que el de la Revolución Industrial}

El juicio de Hassabis sobre el impacto social de la AGI es sumamente serio. Compara la AGI con la Revolución Industrial, pero subraya que tanto su escala como su velocidad superan por mucho a las de aquella.

\begin{importantbox}{Las lecciones históricas de la Revolución Industrial}
Antes de la Revolución Industrial, la mortalidad infantil llegaba a cerca del 40\%. El avance tecnológico que trajo la Revolución Industrial acabó por mejorar enormemente la vida humana, pero \textbf{a la humanidad le tomó un siglo entero digerir sus efectos secundarios}: el trabajo infantil, los barrios marginales, el conflicto de clases, la contaminación ambiental.

El cambio que traerá la AGI será de ``diez veces la escala, diez veces la velocidad''; esto significa que la sociedad tendrá que \textbf{comprimir en diez años} ese proceso de adaptación de cien años. Debemos planear con anticipación, en lugar de remediar las cosas después de que ocurran.
\end{importantbox}

\subsection{El riesgo de la sustitución de la fuerza de trabajo}

Hassabis reconoce que el \textbf{riesgo de sustitución de la fuerza de trabajo} que trae la AGI es real y grave. Presenta una refutación explícita a la postura de tecnooptimistas como Marc Andreessen:

\begin{warningbox}{Refutación a Andreessen: esta vez es distinto}
Marc Andreessen y otros suelen recurrir a analogías históricas para argumentar que el avance tecnológico termina por crear más empleo: ``ninguna revolución tecnológica anterior provocó desempleo masivo, y esta tampoco lo hará''.

La refutación de Hassabis es la siguiente: \textbf{la escala y la velocidad de esta vez hacen que la analogía histórica deje de ser aplicable}. Las revoluciones tecnológicas del pasado (la máquina de vapor, la electricidad, internet) sustituyeron la \textbf{fuerza física} humana o \textbf{capacidades cognitivas específicas}, mientras que la AGI sustituirá la \textbf{capacidad cognitiva integral}, algo sin precedente en la historia humana. Usar patrones del pasado para predecir un suceso sin precedente resulta insuficiente desde el punto de vista lógico.
\end{warningbox}

\subsection{El mecanismo de distribución de la riqueza}

Ante la inquietud de que ``las ganancias de la IA se concentran en unas cuantas empresas'', Hassabis señala un mecanismo importante que ya existe:

\begin{knowledgebox}{El papel de los fondos de pensiones y los fondos soberanos}
Entre los principales inversionistas de las empresas líderes en IA están los \textbf{fondos de pensiones} (pension funds) y los \textbf{fondos soberanos de riqueza} (sovereign wealth funds). Esto significa que la riqueza que genera la IA no se concentra por completo en manos de unos cuantos fundadores e inversionistas de capital de riesgo: una gran cantidad de personas comunes posee de manera indirecta acciones de empresas de IA a través de sus planes de pensión.

Esto no es una solución perfecta, pero ofrece un canal de redistribución de la riqueza que ya existe y que no requiere diseñar una institución enteramente nueva. La verdadera pregunta es: ¿basta esta distribución indirecta para hacer frente al cambio estructural del empleo en la era de la AGI?
\end{knowledgebox}

\subsection{Resumen de la sección}

El impacto socioeconómico de la AGI superará con mucho al de la Revolución Industrial, y el riesgo de sustitución de la fuerza de trabajo no se puede resolver con una simple analogía histórica. Los fondos de pensiones y los fondos soberanos existentes ofrecen un canal parcial de distribución de la riqueza, pero quizá no basten para hacer frente al choque estructural que trae la sustitución cognitiva integral.

%% ==================== Section 7 ====================
